\section*{Assumptions}
\begin{itemize}

    \item \textbf{Block size range.} The invention is demonstrated and evaluated for block sizes $L \in \{256, 512, 1024, 2048, 4096\}$ bits.
    The approach generalizes to any block size; overhead scaling behavior ($O(1/\sqrt{L})$) holds for all $L$.

    \item \textbf{Error rate regime.} The primary operating regime is BER up to approximately 5\% of block bits (i.e., up to $\sim$200 bit flips per 4096-bit block with CRC-32).
    At very high BER ($> 10\%$), success rate degrades gracefully; the scheme does not guarantee correction at arbitrarily high error rates.

    \item \textbf{Availability of a clean baseline.} The coverage-ordered solver requires a set of baseline CRC digests computed over the uncorrupted data.
    These are stored alongside the (possibly corrupted) data block or embedded into the payload, and are assumed to be reliable (e.g., stored in more robust memory, or re-derivable from a checkpointed copy).

    \item \textbf{Key availability for Feistel permutation.} The Feistel permutation requires a per-block key for the permutation.
    This key must be consistently available at both encoding and decoding time.
    For applications without a natural secret key (e.g., non-cryptographic storage), a fixed or publicly known key is sufficient for the burst-equalization benefit.

    \item \textbf{Embedded metadata variant applicability.} The MILP-based embedded metadata variant assumes the data payload consists of structured values whose bit representation has sufficient degrees of freedom to absorb CRC constraint perturbations with acceptable distortion.
    This is satisfied by quantized neural network weights (8-bit or higher), floating-point mantissa bits, and sensor readings represented with headroom.
    The variant is not applicable to data that must be stored at exact bit-for-bit precision.

    \item \textbf{Solver runtime.} The coverage-ordered solver's runtime scales with the number of mismatched nodes and the enumeration depth (maximum flip count attempted per node).
    For BER $\leq 5\%$ at 4096 bits with CRC-32, solver runtime is on the order of seconds on commodity CPU hardware.
    Hardware acceleration (e.g., FPGA or dedicated silicon) can reduce this to microseconds.

    \item \textbf{CRC collision probability.} CRC-$B$ hashes have a false-positive probability of $2^{-B}$ per check.
    CRC-32 provides $< 2.3 \times 10^{-10}$ false-positive probability per hash node, which is negligible for the group sizes considered.

\end{itemize}
